% ===========================================================================
% Shared notation for the internal note and the compendium.
%
% Conventions follow the statistical-physics usage of the diffusion literature
% (Song et al. 2021 for the SDE formulation; Mezard-Montanari 2009 for the
% graphical-model side). Only publicly available sources are cited.
%
% Symbol decisions that matter, fixed here once:
%
%   1. K (or K_theta) is the MARKOV TRANSITION KERNEL, everywhere and only.
%      The number of quadrature grid points is N_g. These were previously W and
%      K respectively, which collided; nothing in either document may now use
%      K for a grid size or W for a kernel.
%
%   2. The chain length is L, and every graphical-model statement is indexed by
%      sites. The ambient dimension equals it, n = L, but the two are written
%      separately because they are different concepts and the reader should not
%      have to infer the identification. The number of observed sequences is M.
%      A translation sentence appears at first use, because much of the
%      diffusion literature uses n for the sample count and d for the dimension.
%
%   5. Vectors are bold lowercase (\vect) and matrices and tensors are bold
%      uppercase (\mat). Scalars, including single sites a_i, stay upright.
%      Anything indexed down to a scalar loses its bold.
%
%   3. The decay factor is written e^{-t} in displayed equations and is never
%      abbreviated to alpha_t.
%
%   4. Gaussian white noise eta is STANDARDISED, <eta(t)eta(t')> = delta(t-t'),
%      so the sqrt(2) is carried explicitly in both the forward and the reverse
%      SDE. This keeps Delta_t = 1 - e^{-2t} consistent throughout.
% ===========================================================================

% --- Diffusion -------------------------------------------------------------
\newcommand{\Pz}{P_0}                       % clean distribution
\newcommand{\Pt}{P_t}                       % noised distribution at time t
\newcommand{\Dt}{\Delta_t}                  % 1 - e^{-2t}
\newcommand{\score}{S}                      % S(x,t) = grad log P_t(x)
\newcommand{\gwn}{\vect{\eta}}              % standardised Gaussian white noise (vector)
\newcommand{\tmin}{t_{\min}}
\newcommand{\tmax}{t_{\max}}

% Posterior mean, in bracket notation rather than E[a|x].
\newcommand{\pmean}[1]{\langle #1 \rangle_{x,t}}
\newcommand{\pmeansite}[1]{\langle a_{#1} \rangle_{x,t}}

% --- Vectors and matrices --------------------------------------------------
% Bold lowercase for vectors, bold uppercase for matrices and tensors.
\newcommand{\vect}[1]{\mathbf{#1}}
\newcommand{\mat}[1]{\mathbf{#1}}
\newcommand{\va}{\vect{a}}                  % clean chain, in R^L
\newcommand{\vx}{\vect{x}}                  % noised chain, in R^L
\newcommand{\vz}{\vect{z}}                  % standard normal draw, in R^L
\newcommand{\Sig}{\mat{\Sigma}}             % covariance
\newcommand{\Id}{\mat{I}}                   % identity

% --- Chain and graphical model --------------------------------------------
\newcommand{\nsites}{L}                     % chain length (ambient dimension n = L)
\newcommand{\nseq}{M}                       % number of observed sequences
\newcommand{\ngrid}{N_{\mathrm{g}}}         % number of quadrature grid points
\newcommand{\gridpt}[1]{u_{#1}}             % grid point
\newcommand{\gridh}{h}                      % grid spacing
\newcommand{\halfwidth}{A}                  % grid half-width, domain [-A, A]
\newcommand{\kernel}{K}                     % transition kernel K(a'|a)
\newcommand{\kernelth}{K_{\params}}         % parameterised transition kernel
\newcommand{\like}{\ell_t}                  % unary likelihood factor
\newcommand{\lmsg}{\overrightarrow{m}}      % forward (left) message
\newcommand{\rmsg}{\overleftarrow{m}}       % backward (right) message
\newcommand{\innov}{\varphi}                % innovation density
\newcommand{\corr}{\rho}                    % autoregression coefficient
\newcommand{\ivar}{q}                       % innovation variance
\newcommand{\ncomp}{C}                      % mixture component count

% --- Estimation ------------------------------------------------------------
\newcommand{\params}{\theta}
\newcommand{\Qfun}{\mathcal{Q}}             % EM auxiliary function
\newcommand{\Xis}{\mat{\Xi}}                % expected transition statistic (a matrix)
\newcommand{\Lik}{\mathcal{L}}              % marginal log-likelihood

% --- Operators -------------------------------------------------------------
\newcommand{\E}{\mathbb{E}}
\newcommand{\Prob}{\mathbb{P}}
\newcommand{\Reals}{\mathbb{R}}
\newcommand{\Cov}{\operatorname{Cov}}
\newcommand{\Var}{\operatorname{Var}}
\newcommand{\KL}{D_{\mathrm{KL}}}
\newcommand{\normal}[2]{\mathcal{N}\!\left(#1,#2\right)}

% --- Estimators ------------------------------------------------------------
\newcommand{\mmse}{\widehat{\va}_{\mathrm{MMSE}}}
\newcommand{\lmmse}{\widehat{\va}_{\mathrm{LMMSE}}}

% --- Pending-data guard -----------------------------------------------------
% Every number in the paper must come from the frozen batch (outputs/frozen/).
% While that batch is running, a slot that still needs a number is marked with
% \needsdata{what}, which renders in loud red with a bracket. It is deliberately
% impossible to miss in a built PDF, so a placeholder cannot be mistaken for a
% result. `tools/check_no_placeholders.sh` greps for it and fails the build if
% any survive.
\providecommand{\needsdata}[1]{{\bfseries\color{red}[PENDING: #1]}}

% --- Repository pointers (shared by the note and the compendium) ------------
\providecommand{\coderef}[2]{\texttt{\small #1}\,$\rightarrow$\,\texttt{\small #2}}
\providecommand{\codefile}[1]{\texttt{\small #1}}
